% =====================================================================
%  MorphBPE — prezentacja (Beamer)
%  Paleta Slayer, kompatybilna z formatką papera
% =====================================================================
\documentclass[aspectratio=169,11pt]{beamer}

% --- Język / kodowanie ---
\usepackage[T1]{fontenc}
\usepackage[utf8]{inputenc}
\usepackage[polish]{babel}
\usepackage{lmodern}

% --- Kolory (paleta Slayer) ---
\usepackage{xcolor}
\definecolor{slayerblue}{RGB}{15, 76, 129}
\definecolor{slayercyan}{RGB}{0, 169, 224}
\definecolor{slayerdark}{RGB}{34, 40, 49}
\definecolor{slayergray}{RGB}{245, 245, 245}
\definecolor{slayergreen}{RGB}{46, 125, 50}
\definecolor{slayeramber}{RGB}{245, 124, 0}
\definecolor{slayerred}{RGB}{198, 40, 40}

% --- Motyw Beamer ---
\usetheme{default}
\usecolortheme{default}
\setbeamertemplate{navigation symbols}{}
\setbeamertemplate{footline}[frame number]
\setbeamerfont{frametitle}{size=\large, series=\bfseries}
\setbeamercolor{frametitle}{fg=slayerblue}
\setbeamercolor{title}{fg=slayerblue}
\setbeamercolor{structure}{fg=slayerblue}
\setbeamercolor{alerted text}{fg=slayerred}
\setbeamercolor{block title}{bg=slayerblue, fg=white}
\setbeamercolor{block body}{bg=slayergray, fg=slayerdark}
\setbeamerfont{block title}{size=\small, series=\bfseries}
\setbeamerfont{block body}{size=\small}

% --- Marginesy ---
\setbeamersize{text margin left=8mm, text margin right=8mm}

% --- Pakiety ---
\usepackage{amsmath, amssymb}
\usepackage{booktabs}
\usepackage{colortbl}
\usepackage{tikz}
\usepackage{pgfplots}
\pgfplotsset{compat=1.18}
\usepackage{tcolorbox}

% --- Metryki ---
\newcommand{\morphed}{\mathrm{MorphED}}
\newcommand{\fert}{\mathrm{Fert}}

% --- Tytuł ---
\title{MorphBPE}
\subtitle{Morfologicznie informowany BPE dla polszczyzny}
\author{Piotr Styła}
\date{\today}

\begin{document}

% =====================================================================
\begin{frame}[plain]
\titlepage
\end{frame}

% =====================================================================
\begin{frame}{Plan prezentacji}
\tableofcontents
\end{frame}

% =====================================================================
\section{Problem}

\begin{frame}{Problem: BPE a fleksja polska}
\begin{block}{Pytanie badawcze}
Czy wprowadzenie morfologicznych ograniczeń do BPE poprawia jakość tokenizacji dla języka polskiego --- intrinsic (alignment) i~extrinsic (LM)?
\end{block}

\bigskip

\textbf{BPE ignoruje strukturę morfologiczną:}
\begin{itemize}
  \item Scalania oparte wyłącznie na częstotliwości par
  \item Tokeny przekraczają granice morfemów
  \item Np.\ \texttt{samorządowych} $\to$ \texttt{sa~|~morządowych} \\
        (rdzeń + sufiks \texttt{-ch} sklejone w~jeden token)
\end{itemize}

\bigskip

\textbf{Konsekwencja dla LM:}
\begin{itemize}
  \item Sufiksy wspólne dla wielu wyrazów nie są rozpoznawane jako wzorzec
  \item Gorsza generalizacja fleksyjna
\end{itemize}
\end{frame}

% =====================================================================
\section{Metoda}

\begin{frame}{MorphBPE: morfologiczne ograniczenie BPE}
\textbf{Idea:} kara za scalenia przekraczające granicę sufiksu.

\bigskip

\textbf{Mechanizm:}
\begin{enumerate}
  \item Morfeusz2 analizuje każde słowo $\to$ granice sufiksów
  \item Przy wyborze pary do scalenia, jeśli scalenie przekracza granicę:
  \[
    \mathrm{score} = \mathrm{freq} \times (1 - \alpha)
  \]
  \item $\alpha = 0$: standardowy BPE (brak kary)
  \item $\alpha = 1$: pełne blokowanie (strict)
  \item $\alpha > 1$: over-penalization (ujemny score)
\end{enumerate}

\bigskip

\textbf{Wariant \texttt{strict\_suffix}:} kara tylko na granice sufiksów \\
(polski = język sufiksalny, prefiksy rzadkie)
\end{frame}

\begin{frame}{Przykład: \texttt{samorządowych}}
\begin{center}
\begin{tabular}{ll}
\toprule
\textbf{BPE} & \texttt{sa~|~morządowych} \\
\textbf{MorphBPE suffix} & \texttt{sa~|~morządowy~|~\textcolor{slayergreen}{ch}} \\
\bottomrule
\end{tabular}
\end{center}

\bigskip

\begin{itemize}
  \item BPE skleja rdzeń z~sufiksem \texttt{-ch} (genitiv/accusativ plural)
  \item MorphBPE zachowuje granicę sufiksu $\to$ LM rozpoznaje wzorzec \texttt{-ch}
  \item Sufiks \texttt{-ch} wspólny dla wielu przypadków $\to$ lepsza generalizacja
\end{itemize}
\end{frame}

\begin{frame}{Korpus i~ewaluacja}
\textbf{Korpus:} 10\,000 przemówień sejmowych (sejm.gov.pl) \\
Split: 80/20 (8000/2000), seed=42

\bigskip

\textbf{Ewaluacja intrinsic:}
\begin{itemize}
  \item \textbf{Fertility} = $\frac{|\mathrm{tokens}|}{|\mathrm{words}|}$ (niżej = bardziej kompaktowo)
  \item \textbf{MorphED} = średnia odległość Levenshteina między segmentacją tokenizatora a~morfologiczną (niżej = lepszy alignment)
  \item Bootstrap CI: 1000 resamples, 95\%
\end{itemize}

\bigskip

\textbf{Ewaluacja extrinsic:}
\begin{itemize}
  \item LSTM (1 warstwa, 64d/128h) i~Transformer (2 warstwy, 4 głowice, 64d)
  \item Cross-entropy i~perplexity na zbiorze testowym
\end{itemize}
\end{frame}

% =====================================================================
\section{Wyniki}

\begin{frame}{Wyniki intrinsic: MorphED}
\begin{center}
\begin{tikzpicture}
  \begin{axis}[
      ybar,
      bar width=14pt,
      width=10cm, height=5cm,
      symbolic x coords={vocab=500, vocab=1000},
      xtick=data,
      ymin=1.5, ymax=2.7,
      ylabel={MorphED (niżej = lepiej)},
      nodes near coords,
      every node near coord/.style={font=\small},
      grid style={gray!20},
      legend style={at={(0.5,-0.28)}, anchor=north, legend columns=2, font=\small},
      enlarge x limits=0.5,
  ]
    \addplot[fill=slayerblue!60, draw=slayerblue] coordinates {
      (vocab=500, 2.503)
      (vocab=1000, 1.838)
    };
    \addlegendentry{BPE}
    \addplot[fill=slayergreen!60, draw=slayergreen] coordinates {
      (vocab=500, 2.414)
      (vocab=1000, 1.712)
    };
    \addlegendentry{MorphBPE suffix}
  \end{axis}
\end{tikzpicture}
\end{center}

\begin{itemize}
  \item vocab=500: $-3{,}6\%$ \quad vocab=1000: $-6{,}8\%$
  \item CI 95\% \textbf{nie pokrywają się} przy obu rozmiarach
  \item Fertility bez zmian ($\Delta < 0{,}1\%$)
\end{itemize}
\end{frame}

\begin{frame}{Wyniki extrinsic: cross-entropy}
\begin{center}
\begin{tikzpicture}
  \begin{axis}[
      ybar,
      bar width=14pt,
      width=10cm, height=5cm,
      symbolic x coords={vocab=500, vocab=1000},
      xtick=data,
      ymin=3.0, ymax=4.0,
      ylabel={Cross-entropy (niżej = lepiej)},
      nodes near coords,
      every node near coord/.style={font=\small},
      grid style={gray!20},
      legend style={at={(0.5,-0.28)}, anchor=north, legend columns=2, font=\small},
      enlarge x limits=0.5,
  ]
    \addplot[fill=slayerblue!60, draw=slayerblue] coordinates {
      (vocab=500, 3.276)
      (vocab=1000, 3.744)
    };
    \addlegendentry{BPE (LSTM)}
    \addplot[fill=slayergreen!60, draw=slayergreen] coordinates {
      (vocab=500, 3.233)
      (vocab=1000, 3.646)
    };
    \addlegendentry{MorphBPE (LSTM)}
  \end{axis}
\end{tikzpicture}
\end{center}

\begin{center}
\small
\begin{tabular}{lrcc}
\toprule
Architektura & Vocab & $\Delta$CE & $\Delta$PPL \\
\midrule
LSTM & 500 & $-1{,}3\%$ & $-4{,}2\%$ \\
LSTM & 1000 & $-2{,}6\%$ & $-9{,}3\%$ \\
Transformer & 500 & $-1{,}3\%$ & $-4{,}0\%$ \\
Transformer & 1000 & $-2{,}7\%$ & $-9{,}0\%$ \\
\bottomrule
\end{tabular}
\end{center}
\end{frame}

\begin{frame}{Analiza błędów per-suffix}
\textbf{Transformer, vocab=1000, $\alpha=0{,}5$} \\
$\Delta < 0$ = MorphBPE lepszy

\bigskip

\begin{columns}[T]
\begin{column}{0.48\textwidth}
\textcolor{slayergreen}{\textbf{MorphBPE zyskuje:}}
\begin{itemize}
  \item \texttt{-ach}: $-11{,}5\%$
  \item \texttt{-ego}: $-3{,}4\%$
  \item \texttt{-ości}: $-3{,}1\%$
  \item \texttt{-enia}: $-1{,}3\%$
\end{itemize}
\smallskip
\textit{Sufiksy złożone --- respektowanie granicy pomaga generalizacji}
\end{column}
\begin{column}{0.48\textwidth}
\textcolor{slayerred}{\textbf{MorphBPE traci:}}
\begin{itemize}
  \item \texttt{-owych}: $+55{,}8\%$
  \item \texttt{-ć}: $+55{,}8\%$
  \item \texttt{-ych}: $+42{,}1\%$
  \item \texttt{-ów}: $+11{,}5\%$
\end{itemize}
\smallskip
\textit{Sufiksy krótkie --- wymuszone oddzielenie = fragmentacja}
\end{column}
\end{columns}

\bigskip
\begin{block}{Wniosek}
Adaptacyjna kara (silna dla długich, słaba dla krótkich sufiksów) mogłaby połączyć korzyści z~mniejszą fragmentacją.
\end{block}
\end{frame}

% =====================================================================
\section{Over-penalization}

\begin{frame}{Over-penalization: $\alpha > 1{,}0$}
\textbf{Odkrycie:} $\alpha = 2{,}0$ daje \textbf{skokową} poprawę nad $\alpha = 1{,}0$

\bigskip

\begin{center}
\begin{tikzpicture}
  \begin{axis}[
      width=11cm, height=5.5cm,
      xlabel={Kara $\alpha$},
      ylabel={Cross-entropy},
      xmin=-0.1, xmax=3.1,
      ymin=3.0, ymax=3.6,
      xtick={0,0.5,1,1.5,2,2.5,3},
      grid=major, grid style={gray!20},
      legend style={at={(0.98,0.98)}, anchor=north east, font=\small},
  ]
    \addplot[slayerblue, mark=*, thick] coordinates {
      (0,3.131) (0.5,3.111) (1,3.090) (1.5,3.092) (2,3.044) (2.5,3.046) (3,3.046)
    };
    \addlegendentry{vocab=500}
    \addplot[slayerred, mark=square*, thick] coordinates {
      (0,3.525) (0.5,3.477) (1,3.432) (1.5,3.420) (2,3.368) (2.5,3.330) (3,3.324)
    };
    \addlegendentry{vocab=1000}
    \draw[dashed, gray] (axis cs:1,3.0) -- (axis cs:1,3.6);
    \node[font=\small, gray] at (axis cs:1,3.55) {over-pen.};
  \end{axis}
\end{tikzpicture}
\end{center}

\begin{itemize}
  \item Skok przy $\alpha=2{,}0$: CE $-1{,}6\%$ (vocab=500), $-1{,}9\%$ (vocab=1000) vs $\alpha=1{,}0$
  \item MorphED również skacze: $2{,}423 \to 2{,}332$ (vocab=500)
  \item Fertility \textbf{spada}: $2{,}449 \to 2{,}388$ --- bardziej konserwatywne segmentacje
\end{itemize}
\end{frame}

\begin{frame}{Dlaczego $\alpha = 2{,}0$ działa?}
\textbf{Mechanizm:}

\begin{itemize}
  \item $\alpha \le 1$: crossing pairs mają score $\ge 0$ --- blokowane tylko gdy są alternatywy
  \item $\alpha > 1$: crossing pairs dostają score $< 0$:
  \begin{itemize}
    \item Gdy wszystkie pary crossing $\to$ najczęstsza wybierana (least negative)
    \item Kolejność scalenia reorganizowana --- merges wewnątrz morfemów faworyzowane
  \end{itemize}
\end{itemize}

\bigskip

\begin{block}{Znaczenie praktyczne}
\begin{itemize}
  \item Mniej tokenów ($-2{,}5\%$ fertility) = tańsze inference LLM
  \item Lepszy LM przy tym samym budżecie (CE $-1{,}6\%$ do $-3{,}1\%$)
  \item Korzyść rośnie ze słownikiem ($-1{,}6\%$ @500 $\to$ $-3{,}1\%$ @1000)
  \item Zerowy koszt dodatkowy --- to ten sam algorytm, inny mnożnik
\end{itemize}
\end{block}
\end{frame}

% =====================================================================
\section{Skalowanie}

\begin{frame}{Skalowanie do produkcji: vocab 8k--32k}
\textbf{W toku} --- trenowanie MorphBPE na pełnym korpusie (8000 tekstów)

\bigskip

\begin{center}
\small
\begin{tabular}{rrrcc}
\toprule
Vocab & $\alpha$ & Fertility & MorphED & Czas \\
\midrule
8000 & 0{,}0 & 1{,}131 & 0{,}574 & 2{,}3h \\
8000 & 1{,}0 & 1{,}141 & 0{,}504 & 1{,}8h \\
\rowcolor{slayergreen!20}
8000 & 2{,}0 & 1{,}149 & \textbf{0{,}456} & 1{,}6h \\
16000 & 0{,}0 & 1{,}051 & 0{,}442 & 2{,}9h \\
16000 & 1{,}0 & 1{,}064 & 0{,}387 & 3{,}2h \\
\bottomrule
\end{tabular}
\end{center}

\bigskip

\begin{itemize}
  \item MorphED spada z~$\alpha=2{,}0$: 0{,}574 $\to$ \textbf{0{,}456} przy vocab=8k
  \item Fertility rośnie minimalnie (+1{,}5\%) --- koszt tokenizacyjny niewielki
  \item Benchmark Dawidek: SI 50{,}8\% $\to$ 58{,}8\%, MFL 2{,}07 $\to$ 1{,}91
  \item Bielik-11B (vocab=32k): fertility=2{,}97 --- MorphBPE vocab=1k: fertility=1{,}84
\end{itemize}
\end{frame}

% =====================================================================
\section{Podsumowanie}

\begin{frame}{Podsumowanie}
\begin{block}{Teza}
MorphBPE \texttt{strict\_suffix} produkuje tokeny lepiej zgrane z~granicami morfemów niż BPE, co przekłada się na niższą cross-entropy LM.
\end{block}

\bigskip

\textbf{Wyniki:}
\begin{itemize}
  \item MorphED $-3{,}6\%$ (vocab=500) do $-6{,}8\%$ (vocab=1000), CI nie pokrywają się
  \item CE LSTM $-1{,}3\%$ do $-2{,}6\%$; Transformer $-1{,}3\%$ do $-2{,}7\%$
  \item PPL $-4{,}0\%$ do $-9{,}3\%$
  \item Fertility bez kosztu ($\Delta < 0{,}1\%$)
\end{itemize}

\bigskip

\textbf{Over-penalization $\alpha=2{,}0$:}
\begin{itemize}
  \item Skokowa poprawa nad $\alpha=1{,}0$: CE $-1{,}6\%$ do $-3{,}1\%$
  \item Fertility \textbf{spada} ($-2{,}5\%$) --- tańsze inference
  \item Optymalna konfiguracja: \texttt{suffix} + $\alpha=2{,}0$
\end{itemize}
\end{frame}

\begin{frame}{Ograniczenia i~przyszłość}
\textbf{Ograniczenia:}
\begin{itemize}
  \item Małe modele LM (LSTM 128, Transformer 64/2L), 500 tekstów treningowych
  \item Wymaga Morfeusz2 --- niedostępny dla wszystkich języków
  \item Traci na sufiksach krótkich (\texttt{-ć}, \texttt{-ych})
  \item Over-penalization testowane tylko na \texttt{suffix}
\end{itemize}

\bigskip

\textbf{Co dalej:}
\begin{itemize}
  \item Skalowanie vocab do 8k/16k/32k (w~toku)
  \item Adaptacyjna kara $\alpha$ zależna od długości sufiksu
  \item Test na innych językach słowiańskich (czeski, rosyjski)
  \item Integracja HuggingFace + eval z~Bielikiem
  \item Większy korpus (Wikipedia/OSCAR)
\end{itemize}
\end{frame}

% =====================================================================
\begin{frame}[plain]
\begin{center}
\vspace{2cm}
{\color{slayerblue}\Huge\bfseries Dziękuję za uwagę\par}
\vspace{1cm}
{\large Pytania?\par}
\vspace{2cm}
{\small\texttt{morphbpe} v0.1.0 --- pip-installable, CLI, 29 testów\par}
{\small\texttt{github.com/.../morphbpe}\par}
\end{center}
\end{frame}

\end{document}
